\begin{textsample}{POS Dim 6 – vem\_ed\_03 – Score 2.00 – vem\_ed\_03/t010.txt}  \label{ex:f6_pos_014}
QUEM É QUEM Stephanie Doscher é diretora do escritório de iniciativas de aprendizagem global da Florida International University (FIU).
Sua formação inclui doutorado em Administração e Supervisão Educacional pela FIU, mestrado em Educação pela Western Washington University e bacharelado em Estudos de História e Teatro pela Emory University (todas as instituições se localizam nos Estados Unidos).
Autora, pesquisadora, \textbf{palestrante} e consultora internacional, seu trabalho recente enfatiza a relação entre diversidade e produção de conhecimento além das fronteiras por meio de intercâmbios virtuais.
Rubricas que desenvolveu para avaliação de desempenho de aprendizagem global de estudantes são usadas em todo o mundo e citadas como parte do quadro de competência global do Programa Internacional de Avaliação de Alunos (PISA), da Organização para a Cooperação e Desenvolvimento Econômico (OCDE).
Coautora de Making Global Learning Universal: Promoting Inclusion and Success for All Students (Stylus \& NAFSA, 2018), um manual para engajar estudantes de graduação nas soluções de problemas globais de forma colaborativa, ela conduz o \textbf{podcast} Making Global Learning Universal e lidera as iniciativas COIL (Collaborative Online International Learning) na FIU.
A FIU vem desenvolvendo COILs (ou PCIs) com as Fatecs desde 2018: já foram 18 projetos colaborativos internacionais nesse período, com as \textbf{unidades} de Americana, Barueri, Bragança Paulista, Campinas, Itapetininga e São Roque.
Por parte da universidade norte-americana, estiveram envolvidos professores de Estudos Asiáticos, Negócios, Educação, História, Hospitalidade e Turismo, Direito, e Estudos sobre Mulheres e Gênero."A primeira parceria institucional da FIU em COIL foi com as Fatecs, e até hoje continua sendo a mais importante", afirma Stephanie.
Sobre o tema que lhe é muito caro, aprendizagem global, Stephanie comenta:"Não é sobre o que você aprende ou onde aprende, é sobre como aprende.
Envolve conectar diferentes perspectivas sobre um tema.
É um processo que engaja pessoas diferentes trabalhando colaborativamente, analisando e resolvendo problemas complexos, que \textbf{transcendem} as fronteiras das diferenças".

% matched lemmas: palestrante, podcast, transcender, unidade
\end{textsample}
